% !TeX root = mean_polynomials_main.tex
% !TeX program = pdflatex
\documentclass{article}
\usepackage{amsmath}
\usepackage{amsfonts}
\usepackage{amssymb}
\setcounter{MaxMatrixCols}{20}

\title{The Cone Generated by Positive Semidefinite Mean Polynomials}
\author{Mehdi Ghasemi, Salma Kuhlmann}
\date{\today}
\usepackage{amsthm}
\usepackage{graphicx}
\newtheorem{theorem}{Theorem}
\newtheorem{proposition}{Proposition}
\newtheorem{definition}{Definition}
\newtheorem{remark}{Remark}
\usepackage{xcolor}
\usepackage{hyperref}

\begin{document}
\maketitle

\begin{abstract}
Certain classes of homogeneous polynomials, such as Positive Semidefinite (PSD) Diagonal-minus-Tail (D-T) 
forms and those derived from the Arithmetic-Geometric Mean inequality—specifically Sums of 
Nonnegative Circuit (SONC) polynomials—can be identified as nonnegative mean polynomials. 
This paper studies PSD forms defined as weighted power means, denoted by $M_{q,p}(Y,w)$, and investigates 
the cone of nonnegative mean real $n$-variate forms of degree $2d$, denoted by $\mathcal{M}_{n,2d}(X,\alpha)$. 
We analyze the relationship between $\mathcal{M}_{n,2d}$ and other fundamental cones in real algebraic geometry: 
the cone of all PSD forms ($\mathcal{P}_{n,2d}$), the cone of Sums of Squares ($\Sigma_{n,2d}$), and the SONC 
cone ($\mathcal{C}_{n,2d}$). Furthermore, we examine the structural properties of the sum of the SOS and SONC 
cones, commonly referred to as the SOSONC cone.
\end{abstract}

\section{Introduction}\label{s01}
The certification of nonnegativity for multivariate polynomials is a central problem in real algebraic geometry 
...with significant applications in polynomial optimization. While deciding whether a general polynomial is 
Positive Semidefinite (PSD) is NP-hard, specific sub-classes of polynomials offer tractable certificates of 
nonnegativity. Notable examples include Sums of Squares (SOS), Positive Semidefinite 
Diagonal-minus-Tail (PSD D-T) forms, and Sums of Nonnegative Circuit (SONC) polynomials, which are derived 
from the Arithmetic-Geometric Mean inequality.

In this paper, we propose a unified framework for these classes by identifying them with
\textit{nonnegative mean polynomials}. We study PSD forms that are constructed as weighted power means,
denoted by $M_{q,p}(Y,w)$, where $Y = (Y_{1},...,Y_{m})$ is a tuple of symbols and $w \in \mathbb{R}_{+}^{m}$ 
is a vector of positive weights.

Specifically, we investigate the \textbf{cone of nonnegative mean real $n$-variate forms of 
degree $2d$}, denoted as $\mathcal{M}_{n,2d}(X,\alpha)$, where $X = (X_{1},...,X_{n})$ and 
$\alpha \in \mathbb{N}_{0}^{n}$. A primary objective of this work is to elucidate 
the relationship between $\mathcal{M}_{n,2d}$ and the following established cones:
\begin{itemize}
    \item $\mathcal{P}_{n,2d}$: The cone of all positive semidefinite (PSD) forms. Note that by definition, 
    $\mathcal{M}_{n,2d} \subseteq \mathcal{P}_{n,2d}$.
    \item $\Sigma_{n,2d}$: The cone of sums of squares (SOS) polynomials.
    \item $\mathcal{C}_{n,2d}$: The cone of sums of nonnegative circuit (SONC) polynomials.
   \item \textbf{SOSONC}: The sum of the SOS and SONC cones, defined by M. Schick as 
   $(\Sigma + \mathcal{C})_{n,2d}$.
\end{itemize}

\subsection*{Structure of the Paper}
The remainder of this paper is organized as follows:

In \S\ref{s02}, we provide the preliminaries on mean polynomials. We introduce the weighted power means 
$M_{q,p}(Y,w)$ for $q, p \in \mathbb{N}_0$, and utilizing Jensen's inequality, we prove that if $p < q$, 
the form $M_{q,p}$ is PSD (Proposition \ref{prop:Monotonicity}). We further demonstrate how standard classes 
of nonnegative polynomials can be recovered as instances of these means.

In \S\ref{s03}, we revisit nonnegative circuit polynomials. We provide necessary and sufficient conditions 
for the nonnegativity of D-T forms and circuit polynomials using the language of mean polynomials, recovering 
established results found in the literature. We formally define the SONC cone $\mathcal{C}_{n,2d}$ 
(Definition~\ref{def:SONC}) and show that the circuit nonnegativity criterion 
(Theorem~\ref{thm:Circuit_PSD}) implies the D-T criterion (Theorem~\ref{thm:DT_PSD}) as a special case; 
the two results are not equivalent.

In \S\ref{s04}, we investigate the relationship between mean polynomials and Sums of Squares. We study 
$M_{2d,p}(X,\alpha)$ with $|\alpha| = 2d$ and $p \mid 2d$. We prove that $M_{2d,0}$ forms are Sums of 
Binomial Squares (SOBS) (Proposition~\ref{prop:M_2d_0_SOBS}). Additionally, for cases $d = 1, 2, 3$, we prove via induction 
and reduction to Hilbert's theorem that these forms are SOS for all $n \in \mathbb{N}$ (Theorem~\ref{thm:M_2d_p_SOS}).

In \S\ref{s05}, we move beyond the SOSONC cone to understand the general form $M_{q,p}(Y,w)$. 
We examine the Choi-Lam form $Q$ and the second Robinson form $\hat{R}$. While $Q$ is known to be 
SONC but not SOS, we establish a representation of $Q$ and $\hat{R}$ as mean polynomials---specifically 
showing that $\hat{R}$ can be expressed as $8M_{1,0}(G, \mathbf{1})$ for a specific linear transformation $G$. 
This identifies $\hat{R}$ as a separating form for the cone of mean polynomials and the SOSONC cone.

The topics of Positivestellensatz, polynomial optimization, and differential-algebra/KKT objectives 
(\S\ref{s06}--\S\ref{s07}) are reserved for a follow-up paper.

\section{Preliminaries on Mean Polynomials}\label{s02}

In this section, we introduce the concept of weighted power means and establishing the fundamental 
inequalities that allow us to construct nonnegative polynomials.

Let $Y = (Y_1, \dots, Y_m)$ be an $m$-tuple of variables (or symbols) and $w = (w_1, \dots, w_m) \in \mathbb{R}_+^m$ 
be a vector of positive weights. We denote the power mean of order $p$ as $M_p(Y, w)$.

\subsection{Jensen's Inequality and Monotonicity}

The nonnegativity of the mean polynomials studied in this paper relies on the monotonicity of the power mean function. 
We prove this using Jensen's inequality.

\begin{proposition}\label{prop:Monotonicity}
If $p < q$, then $M_p(Y, w) \le M_q(Y, w)$ for any non-negative input $Y$.
\end{proposition}

\begin{proof}
Consider the function $\varphi(x) = x^{\frac{2q}{p}}$ with $q > p$. The second derivative of this function is:
\[
\varphi''(x) = \frac{2q}{p} \left( \frac{2q}{p} - 1 \right) x^{\frac{2q}{p} - 2}
\]
Since $q > p$, the coefficient $\frac{2q}{p}(\frac{2q}{p}-1)$ is strictly positive. Consequently, 
$\varphi(x)$ is a convex function on $\mathbb{R}_+$.

By Jensen's inequality, for any convex function $\varphi$ and weights $w_i$, we have:
\[
\varphi\left( \frac{\sum w_i x_i}{\sum w_i} \right) \le \frac{\sum w_i \varphi(x_i)}{\sum w_i}
\]
Substituting $x_i$ with $x_i^p$ in the inequality above, we obtain:
\[
\varphi\left( \frac{\sum w_i x_i^p}{\sum w_i} \right) \le \frac{\sum w_i \varphi(x_i^p)}{\sum w_i}
\]
Substituting the definition $\varphi(t) = t^{q/p}$ (noting that the factor of 2 in the degree is preserved in the 
structure of the resulting forms), this inequality simplifies to:
\[
\left( \frac{\sum w_i x_i^p}{\sum w_i} \right)^{\frac{q}{p}} \le \frac{\sum w_i x_i^q}{\sum w_i}
\]
Taking the $q$-th root, this is equivalent to $M_p(Y, w) \le M_q(Y, w)$.
\end{proof}

\subsection{The Cone of Mean Polynomials}

Using the inequality established above, we can define a class of nonnegative polynomials. By raising both sides of 
the inequality $M_p \le M_q$ to the power $c = \operatorname{lcm}(q, p)$, we eliminate rational exponents. 
We define the \textbf{mean polynomial form} $\mathcal{M}_{q,p}$ as:
\[
\mathcal{M}_{q,p}(Y, w) = M_q^c(Y, w) - M_p^c(Y, w)
\]
By Proposition~\ref{prop:Monotonicity}, $\mathcal{M}_{q,p}$ is a nonnegative polynomial function in variables $Y$ for all $q > p$.

\subsubsection*{Properties of Power Means}
We recall the following basic limiting properties of the mean function $M_p(x, w)$ that are relevant to the boundaries 
of these cones:

\begin{itemize}
    \item \textbf{Minimum:} $\lim_{p \to -\infty} M_p(x, w) = \min(x_1, \dots, x_n)$.
    \item \textbf{Maximum:} $\lim_{p \to \infty} M_p(x, w) = \max(x_1, \dots, x_n)$.
    \item \textbf{Geometric Mean:} $\lim_{p \to 0} M_p(x, w) = \left( \prod_{i=1}^n x_i^{w_i} \right)^{1 / \sum w_i}$.
\end{itemize}

Finally, we note that the equality $M_p(x, w) = M_q(x, w)$ holds if and only if all elements are equal, i.e., 
$x_1 = x_2 = \dots = x_n$.

\section{Nonnegative Diagonal minus Tail and Circuit Polynomials}\label{s03}

In this section, we investigate the conditions for nonnegativity of Diagonal-minus-Tail (D-T) forms and general 
Circuit polynomials. We utilize the language of mean polynomials to re-derive established results 
found in \cite{GM01} and \cite{DIdW01}.

\begin{theorem}\label{thm:DT_PSD}
Let $f(x)$ be a Diagonal-minus-Tail form defined as:
\[
   f(x) = \sum_{i=1}^n \beta_i x_i^{2d} - \mu x^{\alpha}
\]
where:
\begin{itemize}
     \item $x = (x_1, \dots, x_n) \in \mathbb{R}^n$;
     \item coefficients satisfy $\beta_i, \mu > 0$;
     \item exponents satisfy $\alpha = (\alpha_1, \dots, \alpha_n) \in \mathbb{N}_0^n$ and
     \[
         \sum_{i=1}^n \alpha_i = 2d.
     \]
\end{itemize}
The function $f(x)$ is positive semidefinite (PSD) on $\mathbb{R}_+^n$ if and only if:
\[
   \mu^{2d} \prod_{i=1}^n \alpha_i^{\alpha_i} \leq (2d)^{2d} \prod_{i=1}^n \beta_i^{\alpha_i}
\]
\end{theorem}

\begin{proof}
The proof relies on the weighted Arithmetic Mean-Geometric Mean (AM-GM) inequality.

\noindent \textbf{(\(\Leftarrow\))} Assume the condition holds. To show $f(x) \geq 0$, it suffices to show 
$\sum_{i=1}^n \beta_i x_i^{2d} \geq \mu x^{\alpha}$.
Apply the weighted AM-GM inequality to the terms $y_i = \frac{\beta_i}{\alpha_i} x_i^{2d}$ with weights 
$w_i = \alpha_i$ (noting $\sum w_i = 2d$). We obtain:
\[
   \frac{\sum_{i=1}^n \alpha_i \left( \frac{\beta_i}{\alpha_i} x_i^{2d} \right)}{\sum_{i=1}^n \alpha_i} \geq \left( \prod_{i=1}^n \left( \frac{\beta_i}{\alpha_i} x_i^{2d} \right)^{\alpha_i} \right)^{1/2d}
\]
Simplifying the LHS gives $\frac{1}{2d} \sum \beta_i x_i^{2d}$. Rearranging yields:
\[
   \sum_{i=1}^n \beta_i x_i^{2d} \geq 2d \left( \prod_{i=1}^n \left( \frac{\beta_i}{\alpha_i} \right)^{\alpha_i} \right)^{1/2d} x^{\alpha}
\]
By the hypothesis, $\mu \leq 2d \prod (\frac{\beta_i}{\alpha_i})^{\frac{\alpha_i}{2d}}$. 
Thus, $\sum \beta_i x_i^{2d} \geq \mu x^\alpha$.

\noindent \textbf{(\(\Rightarrow\))} Conversely, assume $f(x)$ is PSD. We prove the condition by contradiction. Suppose:
\[
   \mu > 2d \left( \prod_{i=1}^n \left( \frac{\beta_i}{\alpha_i} \right)^{\alpha_i} \right)^{1/2d}
\]
We construct a test point $x^* \in \mathbb{R}_+^n$ where the AM-GM equality holds. 
Let $x_i^* = (\frac{\alpha_i}{\beta_i})^{1/2d}$.
Evaluating the sum term:
\[
   \sum_{i=1}^n \beta_i (x_i^*)^{2d} = \sum_{i=1}^n \beta_i \left(\frac{\alpha_i}{\beta_i}\right) = \sum_{i=1}^n \alpha_i = 2d
\]
Evaluating the monomial term:
\[
   (x^*)^{\alpha} = \prod_{i=1}^n \left( \frac{\alpha_i}{\beta_i} \right)^{\frac{\alpha_i}{2d}} = \left( \prod_{i=1}^n \left( \frac{\alpha_i}{\beta_i} \right)^{\alpha_i} \right)^{1/2d}
\]
Substituting into $f(x^*)$:
\[
   f(x^*) = 2d - \mu \left( \prod_{i=1}^n \left( \frac{\alpha_i}{\beta_i} \right)^{\alpha_i} \right)^{1/2d}
\]
Under the assumption that the condition fails (strictly greater $\mu$), we have $f(x^*) < 0$, contradicting that $f$ is PSD.
\end{proof}

We now generalize this result to circuit polynomials, where the tail term is supported by a convex combination of the 
exponents of the diagonal terms.

\begin{theorem}\label{thm:Circuit_PSD}
Let $\alpha(0), \dots, \alpha(r) \in \mathbb{N}^n$ be exponent vectors. Let $\beta \in \mathbb{N}^n$ be a convex 
combination of the $\alpha(j)$ vectors, such that $\beta = \sum_{j=0}^r \lambda_j \alpha(j)$ with $\lambda_j > 0$ 
and $\sum \lambda_j = 1$.
Consider the polynomial $f(X)$ with coefficients $f_{\alpha(j)} > 0$:
\[
   f(X) = \sum_{j=0}^r f_{\alpha(j)} X^{\alpha(j)} - f_{\beta} X^{\beta}
\]
The polynomial $f(X)$ is PSD for all $X \in \mathbb{R}_+^n$ if and only if:
\[
   |f_{\beta}| \leq \prod_{j=0}^r \left( \frac{f_{\alpha(j)}}{\lambda_j} \right)^{\lambda_j}
\]
\end{theorem}

\begin{proof}
Let $Y_j = \frac{f_{\alpha(j)}}{\lambda_j} X^{\alpha(j)}$ for $j = 0, \dots, r$. By the weighted AM-GM inequality:
\[
   \sum_{j=0}^r \lambda_j Y_j \geq \prod_{j=0}^r Y_j^{\lambda_j}
\]
Substituting $Y_j$:
\[
   \sum_{j=0}^r \lambda_j \left( \frac{f_{\alpha(j)}}{\lambda_j} X^{\alpha(j)} \right) \geq \prod_{j=0}^r \left( \frac{f_{\alpha(j)}}{\lambda_j} X^{\alpha(j)} \right)^{\lambda_j}
\]
The LHS simplifies to $\sum f_{\alpha(j)} X^{\alpha(j)}$. The RHS simplifies to:
\[
   \left( \prod_{j=0}^r \left( \frac{f_{\alpha(j)}}{\lambda_j} \right)^{\lambda_j} \right) X^{\sum \lambda_j \alpha(j)} = \left( \prod_{j=0}^r \left( \frac{f_{\alpha(j)}}{\lambda_j} \right)^{\lambda_j} \right) X^{\beta}
\]
Thus, if $|f_\beta| \le \prod (\frac{f_{\alpha(j)}}{\lambda_j})^{\lambda_j}$, it follows that $\sum f_{\alpha(j)} X^{\alpha(j)} \ge |f_\beta| X^\beta \ge f_\beta X^\beta$, proving nonnegativity. The necessary condition follows by selecting $X$ such that the variance of terms $Y_j$ vanishes.
\end{proof}

\begin{definition}\label{def:SONC}
A polynomial whose coefficients satisfy the condition in Theorem~\ref{thm:Circuit_PSD} is called a 
\textbf{nonnegative circuit polynomial}.
The \textbf{SONC cone} $\mathcal{C}_{n,2d}$ is the conic hull of all nonnegative circuit polynomials of 
degree $2d$ in $n$ variables.
\end{definition}

\begin{remark}\label{rem:SONCinMEANS}
From the proof of the theorem \ref{thm:Circuit_PSD}, it is evident that every nonnegative circuit polynomial 
can be expressed as a mean polynomial of the form $M_{1,0}(Y,w)$, where 
$Y_j = \frac{f_{\alpha(j)}}{\lambda_j} X^{\alpha(j)}$ and $w_j = \lambda_j$.
\end{remark}

\begin{definition}\label{def:SOSONC}
The \textbf{SOSONC cone} $(\Sigma + \mathcal{C})_{n,2d}$ is the Minkowski sum of the SOS cone $\Sigma_{n,2d}$
and the SONC cone $\mathcal{C}_{n,2d}$, as introduced in \cite{DresslerSchick2025}.
\end{definition}

\subsection{Diagonal-minus-Tail vs Circuit Forms}

Theorem~\ref{thm:Circuit_PSD} implies Theorem~\ref{thm:DT_PSD} as a special case; the converse does not 
hold in general, so the two results are \emph{not} equivalent. Every D-T form is a circuit polynomial 
(with a single tail term), but circuit polynomials allow an arbitrary number of diagonal terms whose 
exponents form a simplex, which is a strictly more general setting.

To see the implication, map the D-T form $f(x) = \sum \beta_i x_i^{2d} - \mu x^{\alpha}$ to the general framework:
\begin{itemize}
    \item The diagonal terms correspond to $\alpha(i) = 2d e_i$ with coefficients $f_{\alpha(i)} = \beta_i$.
    \item The tail term exponent is $\beta = \alpha$.
    \item The convex combination weights are $\lambda_i = \frac{\alpha_i}{2d}$. 
    Hence
    \[
        \sum \lambda_i = \frac{1}{2d} \sum \alpha_i = 1.
    \]
\end{itemize}
Applying the condition from Theorem \ref{thm:Circuit_PSD}:
\[
   \mu \leq \prod_{i=1}^n \left( \frac{\beta_i}{\alpha_i/2d} \right)^{\alpha_i/2d} = \left( \prod_{i=1}^n \left( \frac{2d \beta_i}{\alpha_i} \right)^{\alpha_i} \right)^{\frac{1}{2d}}
\]
Raising both sides to the power $2d$ yields $\mu^{2d} \leq (2d)^{2d} \prod (\frac{\beta_i}{\alpha_i})^{\alpha_i}$, 
recovering the condition in Theorem \ref{thm:DT_PSD}.

\section{\texorpdfstring{Investigation of $M_{2d,p}(X,\alpha)$ and Connection to SOBS and SOS}{Investigation of M2d,p(X,alpha) and Connection to SOBS and SOS}}
\label{s04}

In this section, we investigate the containment of the mean polynomial forms $M_{2d,p}(X, \alpha)$ within the 
cones of Sums of Binomial Squares (SOBS) and Sums of Squares (SOS).

We first address the case where $p=0$, which corresponds to the geometric mean.

\begin{proposition}\label{prop:M_2d_0_SOBS}
The form $M_{2d,0}(X, \alpha)$ is a PSD Diagonal-minus-Tail (D-T) form and is a Sum of Binomial Squares (SOBS) 
for any $\alpha \in \mathbb{N}_0^n$ with $|\alpha|=2d$.
\end{proposition}

\begin{proof}
The fact that $M_{2d,0}$ is a PSD D-T form follows from the definitions established in Section \ref{s02}. 
The proof that such forms are SOBS is established in \cite{GM01}.
\end{proof}

We now consider the cases where $p > 0$. We prove that for small degrees $2d \in \{2, 4, 6\}$, these forms are not 
only PSD but also SOS.

\begin{theorem}\label{thm:M_2d_p_SOS}
The forms $M_{2d,p}(X, \alpha)$ are Sums of Squares (SOS) for all $n \in \mathbb{N}$ and any $\alpha \in \mathbb{N}^n$ 
with $|\alpha|=2d$, for the cases:
\begin{itemize}
    \item $d=1$ (Quadratic forms), $p \in \{0, 1\}$;
    \item $d=2$ (Quartic forms), $p \in \{0, 1, 2\}$;
    \item $d=3$ (Sextic forms), $p \in \{0, 1, 2, 3\}$.
\end{itemize}
\end{theorem}

\begin{proof}
We analyze each degree $2d$ separately.

\noindent \textbf{Case $d=1$ (Degree 2):}
The forms $M_{2,0}$ and $M_{2,1}$ are PSD quadratic forms. By Hilbert's 1888 Theorem, every nonnegative quadratic 
form is a sum of squares. Thus, they are SOS.

\noindent \textbf{Case $d=2$ (Degree 4):}
For $p=0$, the form is SOBS (and thus SOS) by Proposition \ref{prop:M_2d_0_SOBS}. We analyze $p=1$ and $p=2$.
\begin{itemize}
    \item \textbf{$p=1$:} The form is $M_{4,1}(X, \alpha) = \frac{1}{4}\sum \alpha_i X_i^4 - (\frac{1}{4}\sum \alpha_i X_i)^4$.
    \begin{enumerate}
        \item For $n=2$ (binary quartics) and $n=3$ (ternary quartics), Hilbert's Theorem states that all PSD forms are SOS. 
        Since $M_{4,1}$ is PSD, it is SOS.
        \item For $n=4$, we consider the uniform weights $\alpha = (1,1,1,1)$.
                 SDP decomposition using Irene \cite{GMIr} confirms $M_{4,1}(X, (1,1,1,1))_{\rm sos}=0$.
               Moreover, an approximate Gram matrix for this instance is given in \cite{GMIr}:
         {\scriptsize
         \setlength\arraycolsep{1.5pt}
         \[
            G = \begin{bmatrix}
            0 & 0 & 0 & 0 & 0 & 0 & 0 & 0 & 0 & 0 & 0 & 0 & 0 & 0 & 0 \\
            0 & 0 & 0 & 0 & 0 & 0 & 0 & 0 & 0 & 0 & 0 & 0 & 0 & 0 & 0 \\
            0 & 0 & 0 & 0 & 0 & 0 & 0 & 0 & 0 & 0 & 0 & 0 & 0 & 0 & 0 \\
            0 & 0 & 0 & 0 & 0 & 0 & 0 & 0 & 0 & 0 & 0 & 0 & 0 & 0 & 0 \\
            0 & 0 & 0 & 0 & 0 & 0 & 0 & 0 & 0 & 0 & 0 & 0 & 0 & 0 & 0 \\
            0 & 0 & 0 & 0 & 0 & 0.25 & -0.01 & -0.04 & -0.01 & -0.04 & -0.04 & -0.01 & -0.04 & -0.04 & -0.04 \\
            0 & 0 & 0 & 0 & 0 & -0.01 & 0.05 & -0.01 & 0.01 & 0.01 & -0.04 & 0.01 & 0.01 & -0.02 & -0.04 \\
            0 & 0 & 0 & 0 & 0 & -0.04 & -0.01 & 0.25 & -0.04 & -0.01 & -0.04 & -0.04 & -0.01 & -0.04 & -0.04 \\
            0 & 0 & 0 & 0 & 0 & -0.01 & 0.01 & -0.04 & 0.05 & 0.01 & -0.01 & 0.01 & -0.02 & 0.01 & -0.04 \\
            0 & 0 & 0 & 0 & 0 & -0.04 & 0.01 & -0.01 & 0.01 & 0.05 & -0.01 & -0.02 & 0.01 & 0.01 & -0.04 \\
            0 & 0 & 0 & 0 & 0 & -0.04 & -0.04 & -0.04 & -0.01 & -0.01 & 0.25 & -0.04 & -0.04 & -0.01 & -0.04 \\
            0 & 0 & 0 & 0 & 0 & -0.01 & 0.01 & -0.04 & 0.01 & -0.02 & -0.04 & 0.05 & 0.01 & 0.01 & -0.01 \\
            0 & 0 & 0 & 0 & 0 & -0.04 & 0.01 & -0.01 & -0.02 & 0.01 & -0.04 & 0.01 & 0.05 & 0.01 & -0.01 \\
            0 & 0 & 0 & 0 & 0 & -0.04 & -0.02 & -0.04 & 0.01 & 0.01 & -0.01 & 0.01 & 0.01 & 0.05 & -0.01 \\
            0 & 0 & 0 & 0 & 0 & -0.04 & -0.04 & -0.04 & -0.04 & -0.04 & -0.04 & -0.01 & -0.01 & -0.01 & 0.25 \\
            \end{bmatrix}
         \]
         }
        For non-uniform weights with 
        $|\alpha|=4$ (e.g., $(2,1,1,0)$), the form effectively involves fewer than 4 variables, reducing to the $n \leq 3$ 
        case, which is SOS.
        \item For $n \geq 5$, any partitions of $4$ into $n$ integers must contain zeros. The form reduces to a case 
        with $n \leq 4$, which we have shown to be SOS.
    \end{enumerate}

    \item \textbf{$p=2$:} The form is $M_{4,2}(X, \alpha) = \frac{1}{4}\sum \alpha_i X_i^4 - (\frac{1}{4}\sum \alpha_i X_i^2)^2$.
    Similar to the $p=1$ case:
    \begin{enumerate}
        \item For $n \leq 3$, it is SOS by Hilbert's Theorem.
        \item For $n=4$ with $\alpha=(1,1,1,1)$, SDP decomposition using Irene \cite{GMIr} confirms $f_{\rm sos}=0$. The following is an approximate Gram matrix \cite{GMIr}:
         {\scriptsize
         \setlength\arraycolsep{1.5pt}
         \[
            G = \begin{bmatrix}
            0 & 0 & 0 & 0 & 0 & 0 & 0 & 0 & 0 & 0 & 0 & 0 & 0 & 0 & 0 \\
            0 & 0 & 0 & 0 & 0 & 0 & 0 & 0 & 0 & 0 & 0 & 0 & 0 & 0 & 0 \\
            0 & 0 & 0 & 0 & 0 & 0 & 0 & 0 & 0 & 0 & 0 & 0 & 0 & 0 & 0 \\
            0 & 0 & 0 & 0 & 0 & 0 & 0 & 0 & 0 & 0 & 0 & 0 & 0 & 0 & 0 \\
            0 & 0 & 0 & 0 & 0 & 0 & 0 & 0 & 0 & 0 & 0 & 0 & 0 & 0 & 0 \\
            0 & 0 & 0 & 0 & 0 & 0.19 & -0 & -0.06 & 0 & -0 & -0.06 & 0 & 0 & 0 & -0.06 \\
            0 & 0 & 0 & 0 & 0 & -0 & 0 & -0 & 0 & -0 & 0 & 0 & 0 & 0 & 0 \\
            0 & 0 & 0 & 0 & 0 & -0.06 & -0 & 0.19 & 0 & -0 & -0.06 & 0 & 0 & 0 & -0.06 \\
            0 & 0 & 0 & 0 & 0 & 0 & 0 & 0 & 0 & -0 & -0 & 0 & 0 & 0 & 0 \\
            0 & 0 & 0 & 0 & 0 & -0 & -0 & -0 & -0 & 0 & 0 & 0 & 0 & 0 & -0 \\
            0 & 0 & 0 & 0 & 0 & -0.06 & 0 & -0.06 & -0 & 0 & 0.19 & 0 & 0 & 0 & -0.06 \\
            0 & 0 & 0 & 0 & 0 & 0 & 0 & 0 & 0 & 0 & 0 & 0 & -0 & -0 & 0 \\
            0 & 0 & 0 & 0 & 0 & 0 & 0 & 0 & 0 & 0 & 0 & -0 & 0 & 0 & 0 \\
            0 & 0 & 0 & 0 & 0 & 0 & 0 & 0 & 0 & 0 & 0 & -0 & 0 & 0 & 0 \\
            0 & 0 & 0 & 0 & 0 & -0.06 & 0 & -0.06 & 0 & -0 & -0.06 & 0 & 0 & 0 & 0.19 \\
            \end{bmatrix}
         \]
         }
        \item For $n \geq 5$, the support of $\alpha$ limits the effective variables to at most 4, reducing to the 
        previous cases.
    \end{enumerate}
\end{itemize}

\noindent \textbf{Case $d=3$ (Degree 6):}
We examine $M_{6,p}(X, \alpha)$ for $p \in \{0, \dots, 3\}$.
\begin{enumerate}
    \item \textbf{$n=2$ (Binary Sextics):} By Hilbert's Theorem, all PSD binary forms of degree 6 are SOS.
    \item \textbf{$n=3$ (Ternary Sextics):}
    \begin{itemize}
        \item For $p=0$, it is SOBS (Prop \ref{prop:M_2d_0_SOBS}).
        \item For $p=1$, we analyze partitions of 6 into 3 parts: $(2,2,2)$, $(3,2,1)$, $(4,1,1)$, $(5,1,0)$, $(6,0,0)$.
        The case $\alpha=(2,2,2)$ is confirmed to be SOS via SDP decomposition using Irene \cite{GMIr}. Cases with
        a zero component reduce to $n \leq 2$ (binary sextics, covered by Hilbert's Theorem). The remaining cases
        $(3,2,1)$ and $(4,1,1)$ are also confirmed to be SOS via SDP decomposition using Irene \cite{GMIr}.
        \item For $p=2$ and $p=3$, SOS membership for all partitions of 6 into 3 non-zero parts is confirmed
        via SDP decomposition using Irene \cite{GMIr}; cases with zero components reduce to $n \leq 2$.
    \end{itemize}
    \item \textbf{$n=4$:}
    Relevant partitions of 6 include $(2,2,1,1)$ and $(3,1,1,1)$. SDP decomposition using Irene \cite{GMIr} confirms these specific instances are SOS.
    Partitions containing zeros reduce to $n \leq 3$.
    \item \textbf{$n \geq 5$:}
    Any partition of 6 into $n \ge 6$ parts (e.g., $(1,1,1,1,1,1)$) or parts with zeros reduces the effective number of 
    variables.
\end{enumerate}
This completes the proof for $d=1, 2, 3$.
\end{proof}

\begin{remark}
Two open questions remain regarding these forms:
\begin{enumerate}
    \item Does Theorem \ref{thm:M_2d_p_SOS} generalize to any degree $d$?
    \item Are these forms contained within the SONC cone for all parameters?
\end{enumerate}
A clean degree-$8$ pilot computed with Irene \cite{GMIr} for $d=4$, $n \in \{3,4\}$,
$p \in \{0,1,2,4\}$, and representative nondegenerate template families yields a
consistent pattern: every tested nondegenerate instance admitted an SDP-based SOS
decomposition, while SONC feasibility held only for the $p \in \{0,4\}$ slices.
For the same families with $p \in \{1,2\}$, the SONC signomial constraints remained
infeasible under local and global solves and under error bounds $10^{-6}$, $10^{-8}$,
and $10^{-10}$. This provides positive evidence for a degree-$8$ extension of
Theorem \ref{thm:M_2d_p_SOS} and negative evidence against a broad SONC
containment statement, while still leaving a theorem-level explanation open.
\end{remark}

\subsection*{Computational Evidence Snapshot (Degrees $8$ and $10$)}
We summarize the clean Irene-based pilot runs for $d \in \{4,5\}$, $n \in \{3,4\}$,
$p \in \{0,1,2,4,5\}$ with template families generated from balanced, mixed, and boundary
weight patterns (one-hot degeneracies tracked separately as inconclusive). The clean records are
reported in:
\begin{itemize}
    \item \path{phase3_runs_clean.jsonl}
    \item \path{phase3_runs_clean_d5.jsonl}
    \item \path{phase3_classification_table.csv}
\end{itemize}

For both $d=4$ and $d=5$, every nondegenerate tested instance admitted an SDP-based SOS
decomposition in our current grid. In contrast, SONC feasibility is support-sensitive: the
nondegenerate families with $p \in \{1,2\}$ that fail at degree $8$ remain infeasible under
local/global SONC solves and tighter tolerance settings, while GP relaxations remain feasible on
all tested nondegenerate families in the same grid. This pattern supports a refined view in which
SOS positivity extends robustly on the tested families, whereas SONC membership requires stronger
structural conditions than the broad conjectural form. Tables~\ref{tab:phase3-classification-d4}
and~\ref{tab:phase3-classification-d5} summarize the refreshed classification output using the
corrected support labels. We retain a placeholder row for $d \geq 6$ in the manuscript tables;
a completed $d=6$ L-C2 SONC/GP slice and an initial $d=6$ L-C1 SOS/SDP batch are tracked
in companion artifacts while the full $d=6$ SOS/SDP extension remains in progress.

% Auto-generated by scripts/phase3_build_manuscript_tables.py

\begin{table}[t]
\centering
\scriptsize
\setlength{\tabcolsep}{4pt}
\caption{Pilot classification summary for $d=4$ (degree $8$). Entries report $(\mathrm{SDP}/\mathrm{SONC}/\mathrm{GP})$ using $\mathrm{S}=$ success, $\mathrm{T}=$ timeout, $\mathrm{F}=$ fail, and $\mathrm{I}=$ inconclusive.}
\label{tab:phase3-classification-d4}
\resizebox{\linewidth}{!}{%
\begin{tabular}{p{1.45cm}p{4.8cm}p{2.15cm}ccccc}
\hline
template & representative $\alpha$ & support class & $n$ values & $p=0$ & $p=1$ & $p=2$ & $p=d$ \\
\hline
uniform & $n=3$: $(3, 3, 2)$; $n=4$: $(2, 2, 2, 2)$ & simplex-like & $\{3, 4\}$ & S/S/S & S/F/S & S/F/S & S/S/S \\
mixed & $n=3$: $(4, 2, 2)$; $n=4$: $(4, 2, 2, 0)$ & simplex-like & $\{3, 4\}$ & S/S/S & S/F/S & S/F/S & S/S/S \\
boundary & $n=3$: $(7, 1, 0)$; $n=4$: $(7, 1, 0, 0)$ & simplex-like & $\{3, 4\}$ & S/S/S & S/F/S & S/F/S & S/S/S \\
sparse & $n=3$: $(8, 0, 0)$; $n=4$: $(8, 0, 0, 0)$ & degenerate & $\{3, 4\}$ & I/I/I & I/I/I & I/I/I & I/I/I \\
\hline
$d\geq 6$ & -- & placeholder & $\{3,4\}$ planned & \multicolumn{1}{c}{deferred} & \multicolumn{1}{c}{deferred} & \multicolumn{1}{c}{deferred} & \multicolumn{1}{c}{deferred} \\
\hline
\end{tabular}
}
\end{table}

\begin{table}[t]
\centering
\scriptsize
\setlength{\tabcolsep}{4pt}
\caption{Pilot classification summary for $d=5$ (degree $10$). Entries report $(\mathrm{SDP}/\mathrm{SONC}/\mathrm{GP})$ using $\mathrm{S}=$ success, $\mathrm{T}=$ timeout, $\mathrm{F}=$ fail, and $\mathrm{I}=$ inconclusive.}
\label{tab:phase3-classification-d5}
\resizebox{\linewidth}{!}{%
\begin{tabular}{p{1.45cm}p{4.8cm}p{2.15cm}ccccc}
\hline
template & representative $\alpha$ & support class & $n$ values & $p=0$ & $p=1$ & $p=2$ & $p=d$ \\
\hline
uniform & $n=3$: $(4, 3, 3)$; $n=4$: $(3, 3, 2, 2)$ & simplex-like & $\{3, 4\}$ & S/S/S & S/F/S & S/F/S & S/S/S \\
mixed & $n=3$: $(5, 3, 2)$; $n=4$: $(5, 3, 2, 0)$ & simplex-like & $\{3, 4\}$ & S/S/S & S/F/S & S/F/S & S/S/S \\
boundary & $n=3$: $(9, 1, 0)$; $n=4$: $(9, 1, 0, 0)$ & simplex-like & $\{3, 4\}$ & S/S/S & S/F/S & S/F/S & S/S/S \\
sparse & $n=3$: $(10, 0, 0)$; $n=4$: $(10, 0, 0, 0)$ & degenerate & $\{3, 4\}$ & I/I/I & I/I/I & I/I/I & I/I/I \\
\hline
$d\geq 6$ & -- & placeholder & $\{3,4\}$ planned & \multicolumn{1}{c}{deferred} & \multicolumn{1}{c}{deferred} & \multicolumn{1}{c}{deferred} & \multicolumn{1}{c}{deferred} \\
\hline
\end{tabular}
}
\end{table}

\paragraph{Deferred higher-degree slice.}
The $d\geq 6$ rows are retained as placeholders for a future extension of the same protocol. No higher-degree computations are included in the present manuscript snapshot.


\subsection*{Claim-Status Synthesis (Frozen $d=4,d=5$ Slice)}
Using the frozen clean artifacts and refreshed classification outputs
(listed in the reproducibility note below), we classify the current
computational evidence as follows.
\begin{itemize}
    \item \textbf{Confirmed (within tested grid).}
    For nondegenerate tested families at $d=4,5$ with $n \in \{3,4\}$, SDP/SOS and GP
    outcomes are consistently feasible in the current runs.
    \item \textbf{Weakened (scope-limited).}
    Any broad SONC containment formulation for the same families is not supported
    uniformly; the evidence favors a support-sensitive SONC statement rather than a
    degree-only claim.
    \item \textbf{Negative (robust failures).}
    For nondegenerate families with $p \in \{1,2\}$, SONC infeasibility persists under
    local/global solves and tighter tolerances in both degree slices, while SDP/SOS and
    GP remain feasible.
\end{itemize}
This synthesis is tied to the classification tables in
Tables~\ref{tab:phase3-classification-d4} and~\ref{tab:phase3-classification-d5},
with a partial higher-degree extension from the completed $d=6$ L-C2 SONC/GP slice.

\paragraph{Reproducibility note.}
The tables above are derived from the frozen $d=4,d=5$ artifacts:
\begin{itemize}
    \item \path{phase3_runs_clean.jsonl}
    \item \path{phase3_runs_clean_d5.jsonl}
    \item \path{phase3_classification_table.csv}
    \item \path{phase3_classification_table.md}
\end{itemize}
A partial higher-degree extension is available for $d=6$.
The completed L-C2 SONC/GP slice and an initial L-C1 SOS/SDP batch are:
\begin{itemize}
    \item \path{phase3_runs_clean_d6_lc2.jsonl}
    \item \path{phase3_pilot_summary_clean_d6_lc2.md}
    \item \path{phase3_sonc_diagnostics_d6_lc2.jsonl}
    \item \path{phase3_runs_clean_d6_lc1_batch1.jsonl}
    \item \path{phase3_pilot_summary_clean_d6_lc1_batch1.md}
\end{itemize}
The manuscript snippet \path{results/phase3_classification_tables.tex} is generated
from the refreshed CSV by the table-builder script in \path{scripts/}.
This keeps updates to the frozen $d=4,d=5$ slice and staged $d \geq 6$
extensions reproducible without manual table editing.

\paragraph{Follow-up evidence (in progress).}
Sparse-family counterexample evidence (CX-1), cross-degree stress tests (CX-2), and
the remaining $d=6$ SOS/SDP extension work are under development and will expand the
claim-status synthesis upon completion.

\section{Beyond the SOSONC Cone}
\label{s05}

In this section, we move towards a comprehensive understanding of the general form $M_{q,p}(Y,w)$ beyond the established 
SOSONC cone. We examine specific classical forms—the Choi-Lam form $Q$ and the Robinson forms $R$ and $\hat{R}$—and 
establish their representation as mean polynomials.

\subsection{Specific Examples}
We consider the following forms in $\mathbb{R}[X,Y,Z,W]$:
\begin{itemize}
    \item The Choi-Lam form $Q$:
    \[
        Q(X,Y,Z,W) = X^4 + Y^4 + Z^4 + W^4 - 4XYZW
    \]
    \item The Robinson form $R$:
    \[
        \begin{aligned}
            R(X,Y,Z) ={}& X^6 + Y^6 + Z^6 \\
            &- (X^4 Y^2 + X^4 Z^2 + X^2 Y^4 + X^2 Z^4 + Y^4 Z^2 + Y^2 Z^4) \\
            &+ 3X^2 Y^2 Z^2
        \end{aligned}
    \]
    \item The second Robinson form $\hat{R}$:
    \[
        \begin{aligned}
            \hat{R}(X,Y,Z,W) ={}& X^2(X-W)^2 + Y^2(Y-W)^2 + Z^2(Z-W)^2 \\
            &+ 2XYZ(X+Y+Z-2W)
        \end{aligned}
    \]
\end{itemize}

It is observed in \cite{MSc01} that $Q$ is a Sum of Nonnegative Circuit (SONC) polynomial. However, using the 
Newton polytope method, one can prove that $Q \in \mathcal{P}_{4,4} \setminus \Sigma_{4,4}$ (it is PSD but not SOS).

Remarkably, there exists a linear relationship between $Q$ and $\hat{R}$ as noted by Reznick \cite{Rez07}:
\begin{equation}\label{eq:R_hat_Q_relation}
    \hat{R}(X-W, Y-W, Z-W, X+Y+Z-W) = 2Q(X,Y,Z,W)
\end{equation}
Since the PSD and SOS cones are closed under linear transformations, the identity above shows that $\hat{R}$ is PSD 
(respectively SOS) if and only if $Q$ is PSD (respectively SOS). Consequently, $\hat{R}$ also serves as an example of a 
form in $\mathcal{P}_{4,4} \setminus \Sigma_{4,4}$.

\subsection{Mean Polynomial Representations}

We now establish the representation of these forms as instances of mean polynomials $M_{q,p}$.

\begin{proposition}\label{prop:ChoiLamRobinson}
The Choi-Lam form $Q$ and the Robinson form $\hat{R}$ can be expressed as mean polynomials of order $1$ and $0$ 
(arithmetic minus geometric means).
\end{proposition}

\begin{proof}
For the Choi-Lam form $Q$, we observe:
\[
    Q(X, Y, Z, W) = 4 \left( \frac{X^4 + Y^4 + Z^4 + W^4}{4} - \sqrt[4]{X^4 Y^4 Z^4 W^4} \right)
\]
This corresponds exactly to the mean form definition with variables being the fourth powers and uniform weights:
\[
    Q(X, Y, Z, W) = 4 M_{1,0}((X^4, Y^4, Z^4, W^4), (1, 1, 1, 1))
\]
For the Robinson form $\hat{R}$, applying the inverse of the linear transformation in \eqref{eq:R_hat_Q_relation}, 
we derive:
\[
    \hat{R}(X,Y,Z,W) = 8 M_{1,0}((\hat{X}^4, \hat{Y}^4, \hat{Z}^4, \hat{W}^4), (1,1,1,1))
\]
where 
\[
    \begin{aligned}
        (\hat{X}, \hat{Y}, \hat{Z}, \hat{W}) = \frac{1}{2}(&X+W-Y-Z, \; Y+W-X-Z, \\
        &Z+W-X-Y, \; W-(X+Y+Z)).
    \end{aligned}
\]
\end{proof}

This representation is significant for distinguishing the cone of mean polynomials from the SOSONC cone. While $Q$ is 
SONC (and thus in the SOSONC cone) because it is a mean of monomials in the standard basis, $\hat{R}$ is a mean of 
\emph{linear forms}. Since the SONC property is not invariant under general linear coordinate changes, $\hat{R}$ 
serves as a separating form: it belongs to the cone generated by generalized mean polynomials (means of linear forms) 
but lies outside the standard SOSONC cone (see \cite[page 31 and Lemma 4.2.5, page 107]{MSc01}).


\section{\texorpdfstring{Positivestellensatz and relation to the cone of $2d$ powers}{Positivestellensatz and relation to the cone of 2d powers}}
\label{s06}

\emph{This section is reserved for a follow-up paper.}

\section{Application to polynomial optimization}
\label{s07}

\emph{This section is reserved for a follow-up paper.}

\begin{thebibliography}{99}

\bibitem{DIdW01} M. Dressler, S. Iliman, and T. de Wolff, An Approach to Constrained Polynomial Optimization via Nonnegative Circuit Polynomials and Geometric Programming, \emph{Journal of Symbolic Computation} (2019) \textbf{91}, 149-172

\bibitem{DIdW02} M. Dressler, S. Iliman, and T. de Wolff. \emph{A Positivstellensatz for Sums of Nonnegative Circuit Polynomials}. SIAM J. Appl. Algebra Geom., 1(1):536--555, 2017.

\bibitem{DresslerSchick2025} M. Dressler, S. Kuhlmann and Moritz Schick, \emph{Study of the cone of sums of squares plus sums of nonnegative circuit forms} Advances in Geometry, 25(1), 127--146 (2025)

\bibitem{GLM} M. Ghasemi, J.B. Lasserre, and M. Marshall, \emph{Lower bounds on the global minimum of a polynomial}, Comput. Optim. Appl. 57 (2014), no. 2, 387--402.

\bibitem{GM01} M. Ghasemi and M. Marshall, \emph{Lower bounds for polynomials using geometric programming}, SIAM J. Optim. 22 (2012), no. 2, 460--473.

\bibitem{GM02} M. Ghasemi and M. Marshall, \emph{Lower bounds for a polynomial on a basic closed semialgebraic set using geometric programming}, 2013, Preprint, arxiv:1311.3726.

\bibitem{GM03} M. Ghasemi, M. Marshall,
\emph{Lower bound for a polynomial on a product of hyper\-ellipsoids using geometric programming},
\newblock Preprint, 2015.
\newblock arXiv:2510.03105.

\bibitem{GMIr} M. Ghasemi, \emph{Irene: A Python package for global optimization problems}, 2025, \url{https://github.com/mghasemi/Irene}.

\bibitem{IdW} Sadik Iliman and Timo De Wolff, \emph{Lower Bounds for Polynomials with Simplex Newton Polytopes Based on Geometric Programming,} SIAM Journal on Optimization 26, no. 2 (January 2016): 1128--46, \url{https://doi.org/10.1137/140962425}

\bibitem{JBL01} J. B. Lasserre. \emph{Global optimization with polynomials and the problem of moments.} SIAM J. Optim., 11(3):796--817, 2001.

\bibitem{Rez07} B. Reznick.
\emph{On Hilbert's construction of positive poly\-nomials}.
\newblock Preprint, 2007.
\newblock arXiv:0707.2156.

\bibitem{MSc01} Moritz Schick \emph{Sums of Squares plus Sums of Nonnegative Circuit Polynomials}, PhD. Dissertation. 2025.

\end{thebibliography}

\bigskip

\noindent\textbf{Mehdi Ghasemi:} mehdi.ghasemi@usask.ca\\
Edmonton Police Service

\noindent\textbf{Salma Kuhlmann:} salma.kuhlmann@uni-konstanz.de\\
University of Konstanz

\end{document}
